\documentclass[12pt, a4paper, listof=totoc, bibliography=totoc, numbers=noenddot, ngerman, headsepline, oneside]{scrbook}
\usepackage{amsmath}
\usepackage[T1]{fontenc}
\usepackage{float}
\usepackage[utf8]{inputenc}
\usepackage[ngerman]{babel}
\usepackage{url}
\usepackage{graphicx} 
\usepackage{pdfpages} 
\usepackage[a4paper, margin=1in]{geometry}
\usepackage[right]{eurosym} %Euro-Zeichen
\usepackage{amssymb}
\usepackage{subfig}
\usepackage{cite} %Quellenangaben
\usepackage{setspace} % Zeilenabstand
\usepackage[ 
   colorlinks,       
   linkcolor=black,   % Farbe interner Verweise 
   filecolor=black,   % Farbe externer Verweise 
   citecolor=black,   % Farbe von Zitaten 
   urlcolor=blue	  % Farbe von Links
   ]{hyperref} %Verlinkungen
\usepackage[figure]{hypcap}
\usepackage[ngerman]{translator}
\usepackage{blindtext} % Lorem-Ipsum-Plugin
\usepackage[acronym, nonumberlist]{glossaries} %% use after hyperref %Glossar-Paket laden
%\usepackage[
	%nonumberlist, %keine Seitenzahlen anzeigen
	%acronym,      %ein Abkürzungsverzeichnis erstellen
	%toc,          %Einträge im Inhaltsverzeichnis
	%section      %im Inhaltsverzeichnis auf section-Ebene erscheinen
	%]
%{glossaries}

\usepackage{listings,xcolor} %Codeanzeige
\usepackage{scrhack}
\usepackage[normalem]{ulem}
\useunder{\uline}{\ul}{}
\usepackage{wrapfig}

\usepackage{makecell}
\usepackage{comment}
\renewcommand{\familydefault}{\sfdefault} %%SCHRIFTART
\usepackage{pifont}

\graphicspath{ {./bilder/} }


\usepackage{chngcntr}
\counterwithout{figure}{chapter}
\counterwithout{table}{chapter}

\definecolor{dkgreen}{rgb}{0,.6,0}
\definecolor{dkblue}{rgb}{0.337, 0.612, 0.839}
\definecolor{dkyellow}{rgb}{204, 255, 0}
\definecolor{lila}{rgb}{0.847, 0.624, 0.855}
\definecolor{hintergrund}{rgb}{0.118, 0.118, 0.118}
\definecolor{blau1}{rgb}{0.573, 0.820, 0.925}
\definecolor{orange}{rgb}{0.839, 0.616, 0.522}
\definecolor{psrot}{rgb}{0.66, 0.18, 0.00}
\definecolor{psblau}{rgb}{0, 0, 0.545}
\definecolor{psbg}{rgb}{0.9176, 0.9490, 0.9804}
\definecolor{psbg2}{rgb}{0.949, 0.949, 0.949}
\definecolor{psbg3}{rgb}{1.0, 0.9843, 0.9608}



\lstset{
    numbers=left, 
    numberstyle=\tiny\color{black}, 
    numbersep=5pt,
    breaklines=true,
    frame=lr,
    escapeinside={(*@}{@*)}, %nicht anzuzeigende Ausdrücke, z.B. für Labels
    language=[Sharp]C,
    showstringspaces=false,
    captionpos=b,
    backgroundcolor=\color{hintergrund},
    basicstyle=\ttfamily\fontsize{9}{10}\selectfont\color{white},
    keywordstyle    = \color{dkblue},
    stringstyle     = \color{orange},
    identifierstyle = \color{blau1},
    commentstyle    = \color{dkyellow},
    emph            =[1]{var},
    emphstyle       =[1]\color{dkblue},
    emph            =[2]{if,and,or,else, return},
    emphstyle       =[2]\color{lila}
    } 
\lstdefinestyle{cli}{
    basicstyle=\ttfamily\fontsize{9}{10}\selectfont\color{white},
    identifierstyle = \color{white},
    numbers=none,
    language=sh,
    stringstyle=\color{dkblue},
    emph            =[1]{dotnet},
    emphstyle       =[1]\color{dkyellow},
}
\lstdefinestyle{tiny}{
    basicstyle=\ttfamily\fontsize{7}{8}\selectfont\color{white},
    numbers=left, 
    numberstyle=\tiny\color{black}, 
    numbersep=5pt,
    breaklines=true,
    frame=lr,
    escapeinside={(*@}{@*)}, %nicht anzuzeigende Ausdrücke, z.B. für Labels
    language=[Sharp]C,
    showstringspaces=false,
    captionpos=t,
    backgroundcolor=\color{hintergrund},
    keywordstyle    = \color{dkblue},
    stringstyle     = \color{orange},
    identifierstyle = \color{blau1},
    commentstyle    = \color{dkyellow},
    emph            =[1]{var},
    emphstyle       =[1]\color{dkblue},
    emph            =[2]{if,and,or,else, return},
    emphstyle       =[2]\color{lila}
}
\lstdefinestyle{ps}{
    numbers=left, 
    numberstyle=\tiny\color{black}, 
    numbersep=5pt,
    breaklines=true,
    frame=none,
    escapeinside={(*@}{@*)}, %nicht anzuzeigende Ausdrücke, z.B. für Labels
    language=sh,
    showstringspaces=false,
    captionpos=b,
    backgroundcolor=\color{psbg3},
    basicstyle=\ttfamily\fontsize{9}{10}\selectfont\color{black},
    keywordstyle    = \color{psblau},
    stringstyle     = \color{psrot},
    identifierstyle = \color{psrot},
    commentstyle    = \color{dkyellow},
    emph            =[1]{var},
    emphstyle       =[1]\color{dkblue},
    emph            =[2]{if,and,or,else, return},
    emphstyle       =[2]\color{lila}
}
    

\renewcommand\lstlistingname{Codeausschnitt}
\renewcommand\lstlistlistingname{Codeverzeichnis}

\geometry{verbose,tmargin=2cm,bmargin=2cm,lmargin=2cm,rmargin=2.3cm} 

\clubpenalty = 10000 \widowpenalty = 10000 \displaywidowpenalty = 10000 

\newcommand{\footfigref}[1]{\footnote{Abb. \ref{#1} auf Seite \pageref{#1}}}

%% Bei Referenzen im Text wird jetzt bei allen Ebenen "Kapitel" vorgestellt, z.b. Kapitel 2, Kapitel 2.2, Kapitel 6.3.2
\addto\extrasngerman{%
    \def\sectionautorefname{Kapitel}%
    \def\subsectionautorefname{Kapitel}%
    \def\subsubsectionautorefname{Kapitel}%
    }

% Vertikaler Abstand zwischen Ende Textblock - Ende Fußzeile --> Abstand der Seitenzahl von Rand erhöhen 
\setlength{\footskip}{10mm}

\RedeclareSectionCommand[%
    beforeskip=0.5\baselineskip,
    afterskip=0.5\baselineskip
]{chapter}

\RedeclareSectionCommand[%
    beforeskip=0.5\baselineskip,
    afterskip=0.5\baselineskip
]{section}

\RedeclareSectionCommand[%
    beforeskip=0.1\baselineskip,
    afterskip=0.1\baselineskip
]{subsection}

\RedeclareSectionCommand[%
    beforeskip=0.1\baselineskip,
    afterskip=0.1\baselineskip
]{subsubsection}

\RedeclareSectionCommand[%
    beforeskip=0.01\baselineskip,
    %%afterskip=0.2\baselineskip
]{paragraph}

\setlength{\abovecaptionskip}{4pt}  % 1pc=12pt 
\setlength{\belowcaptionskip}{0pt}
%\setlength{\textfloatsep}{4pt}
\setlength{\intextsep}{1pc}

%% Verkleinerung der Textgröße unter Abbildungen
\addtokomafont{caption}{\small}

%%falsche Default-Silbentrennung überschreiben
\hyphenation{Soft-ware-ent-wick-lung}


\renewcommand*{\glspostdescription}{}

%Glossar-Befehle anschalten
%\makeglossaries
 
\KOMAoptions{parskip=full*}

% ändert Titelschriftart in Serifen-Normalschriftart
\addtokomafont{disposition}{\rmfamily} 

\makenoidxglossaries

\loadglsentries{glossar.tex}

\newcommand{\type}{Software Engineering I - Paper}
\newcommand{\topic}{Grow Green}
\newcommand{\subtopic}{eine ??? Webanwendung}
\newcommand{\Alex}{Alexander Betke (77203378972)}
\newcommand{\Maja}{Maja Günther (72205653284)}
\newcommand{\Theo}{Theo Leuthardt (77205844868)}
\newcommand{\Josh}{Josh Nicolai Tischer (77207542779)}
\newcommand{\Domi}{Domenik Wilhelm (77207300494)}
\newcommand{\company}{Freie Universität Berlin}
\newcommand{\jahrgang}{2023}
\newcommand{\fachbereich}{Duales Studium Wirtschaft · Technik}
\newcommand{\studiengang}{Informatik}
\newcommand{\modul}{IT2101 - Software Engineering I}
\newcommand{\betreuerHS}{Benjamin Kretzmer}
\newcommand{\wordCount}{2800}


\begin{document}
\author{}

\title{
\normalfont\endgraf\rule{\textwidth}{1pt}
\begingroup
	\centering
	\linespread{1.5}
	\huge\topic
\endgroup
\linespread{1.5}
\ \\ % Falls kein Subtopic, auskommentieren
\large\subtopic % Falls kein Subtopic, auskommentieren
\normalfont\endgraf\rule{\textwidth}{1pt}
}

\date{\large vorgelegt am 05. November 2024}

\publishers{
	\begin{tabular}{l l}
    \textbf{\normalsize{}} & \normalsize{}  \tabularnewline
    \textbf{\normalsize{}} & \normalsize{}  \tabularnewline
	\textbf{\normalsize{Modul:}} & \normalsize{\modul}  \tabularnewline
	\textbf{\normalsize{Namen:}} & \normalsize{\Alex}  \tabularnewline
 	\textbf{\normalsize{}} & \normalsize{\Maja}  \tabularnewline
   	\textbf{\normalsize{}} & \normalsize{\Theo}  \tabularnewline
    \textbf{\normalsize{}} & \normalsize{\Josh}  \tabularnewline
    \textbf{\normalsize{}} & \normalsize{\Domi}  \tabularnewline
	\textbf{\normalsize{Studienjahrgang:}} & \normalsize{\jahrgang}  \tabularnewline
	\textbf{\normalsize{Fachbereich:}} & \normalsize{\fachbereich} \tabularnewline
	\textbf{\normalsize{Studiengang:}} & \normalsize{\studiengang} \tabularnewline
	\textbf{\normalsize{Betreuerin Hochschule:}} & \normalsize{\betreuerHS} \tabularnewline
    \textbf{\normalsize{Anzahl der Wörter:}} & \normalsize{\wordCount} \tabularnewline
    \tabularnewline
    \tabularnewline
	\end{tabular}
    \begin{tabular}{p{4.6em} p{0.3em} p{4.6em} p{0.3em} p{4.6em} p{0.3em} p{4.6em} p{0.3em} p{4.6em}}
        \hspace{4.5cm} && \hspace{4.5cm} && \hspace{4.5cm} && \hspace{4.5cm} && \hspace{4.5cm}\\
        \cline{1-1}\cline{3-3}\cline{5-5}\cline{7-7}\cline{9-9} 
        \small{Alexander Betke} && \small{Maja Günther} && \small{Theo Leuthardt} && \small{Josh Tischer} && \small{Domenik Wilhelm}
    \end{tabular}
	}

\titlehead{\begin{center}
    \includegraphics[height=0.031\textheight]{bilder/HWR.png}
    \hfill
    \includegraphics[height=0.031\textheight]{bilder/Polizei.png}
    \hfill
    \includegraphics[height=0.023\textheight]{bilder/FUBIT.png}
    \hfill
    \includegraphics[height=0.026\textheight]{bilder/BDR.jpg}
    \hfill
    \includegraphics[height=0.029\textheight]{bilder/Alstom.png}
    \end{center}
    }

\maketitle
\onehalfspacing 
\pagenumbering{Roman}
\chapter*{Abstract}
\addcontentsline{toc}{chapter}{Abstract}

\clearpage

\newpage

\tableofcontents{}
\addcontentsline{toc}{chapter}{Inhaltsverzeichnis}

\clearpage

\addcontentsline{toc}{chapter}{Akronyme}
\printnoidxglossary[type=\acronymtype]
\clearpage

\addcontentsline{toc}{chapter}{Glossar}
\printnoidxglossary
\clearpage

\chapter{Einleitung}\label{sec:intro}
Kurz Bsp. \cite{Zen24}
\pagenumbering{arabic}
Kurz Bsp. \cite{Zen24}
\chapter{Rahmenbedingungen}\label{sec:rahmen}
\section{Organisatorisch}
 
\section{Technisch}

\section{Rechtlich}

\chapter{Anforderungsanalyse}\label{sec:anforderungen}

\chapter{Verwendete Software und Tools}\label{sec:technik}

\clearpage
\section{Evaluation der erreichten Ergebnisse}

\chapter{Fazit und Ausblick}\label{sec:fazit}

\clearpage
\pagenumbering{Alph}

\bibliography{literatur}

\bibliographystyle{hwrbib}


\listoffigures
\lstlistoflistings

\listoftables
\addchap{Anhang}


\chapter*{Ehrenwörtliche Erklärungen}
\addcontentsline{toc}{chapter}{Ehrenwörtliche Erklärung}


Wir erklären ehrenwörtlich:
\begin{enumerate}
	\item dass wir unser Paper selbstständig verfasst habe,
	\item dass wir die Übernahme wörtlicher Zitate aus der Literatur sowie die Verwendung der Gedanken anderer Autoren an den entsprechenden Stellen innerhalb der Arbeit gekennzeichnet habe,
	\item dass wir unser Paper bei keiner anderen Prüfung vorgelegt haben.
\end{enumerate}
Wir sind uns bewusst, dass eine falsche Erklärung rechtliche Folgen haben wird.\\
\\ \\
\begin{tabular}{lp{2em}l} 
 \hspace{4cm}   && \hspace{7cm} \\\cline{1-1}\cline{3-3} 
 \small{Ort, Datum} && \small{Alexander Betke} \\
 \\\\
 \hspace{4cm}   && \hspace{7cm} \\\cline{1-1}\cline{3-3} 
 \small{Ort, Datum} && \small{Maja Günther} \\
 \\\\
  \hspace{4cm}   && \hspace{7cm} \\\cline{1-1}\cline{3-3} 
 \small{Ort, Datum} && \small{Theo Leuthardt} \\
 \\\\
  \hspace{4cm}   && \hspace{7cm} \\\cline{1-1}\cline{3-3} 
 \small{Ort, Datum} && \small{Josh Nicolai Tischer} \\
 \\\\
  \hspace{4cm}   && \hspace{7cm} \\\cline{1-1}\cline{3-3} 
 \small{Ort, Datum} && \small{Domenik Wilhelm} \\
\end{tabular}

\end{document}
